%! AP Statistics — Unit 4, Lesson 4.7 — Group Activity (Answer Key)
\documentclass[10pt]{article}
\usepackage{apstats-article}
\usepackage{apstats-key}

\begin{document}

\pageheader{Unit 4, Lesson 4.7}{Group Activity: Building and Reading Distributions --- Answer Key}
\namepartnerperiod

\vspace{0.1in}

\begin{tcolorbox}[colback=sky!60, colframe=navy, title=\textbf{Tier R --- Remediate},
    arc=2mm, boxrule=0.8pt, left=4mm, right=4mm, top=3mm, bottom=3mm]
\small
A survey of 50 students asked how many siblings each student has. Results:
\begin{center}
\renewcommand{\arraystretch}{1.4}
\begin{tabular}{@{}cccc@{}}
\textbf{Siblings ($x$)} & \textbf{Frequency} & \textbf{$P(X=x)$} & \textbf{$P(X \le x)$} \\
\hline
0 & 8  & \ans{0.16} & \ans{0.16} \\
1 & 20 & \ans{0.40} & \ans{0.56} \\
2 & 15 & \ans{0.30} & \ans{0.86} \\
3 & 7  & \ans{0.14} & \ans{1.00} \\
\end{tabular}
\end{center}

\begin{enumerate}[leftmargin=1.5em, itemsep=8pt]
  \item Fill in the $P(X=x)$ column. Divide each frequency by the total.

        \vspace{0.1in}

  \item Do these probabilities form a valid distribution? Check both conditions.\\
        \ansline{Valid: every value is $\ge 0$ and $0.16+0.40+0.30+0.14 = 1.00$.}

  \item Fill in the $P(X \le x)$ column.

        \vspace{0.05in}

  \item Using your table, find $P(X \ge 2)$. Show your work.

        \noindent\begin{minipage}[t][0.3in][t]{\linewidth}\ans{$P(X \ge 2) = 0.30 + 0.14 = 0.44$ \quad (or $1 - 0.56 = 0.44$)}\end{minipage}

  \item Is it more likely that a randomly chosen student has 0 siblings or 3 siblings?
        Explain.\\
        \ansline{0 siblings: $P(X=0) = 0.16 > P(X=3) = 0.14$.}
\end{enumerate}
\end{tcolorbox}

\vspace{0.12in}

\begin{tcolorbox}[colback=goldbg!60, colframe=goldacc, title=\textbf{Tier A --- Approaching Proficiency},
    arc=2mm, boxrule=0.8pt, left=4mm, right=4mm, top=3mm, bottom=3mm]
\small
$X$ = number of books a student reads during summer break. The probability distribution is:
\begin{center}
\renewcommand{\arraystretch}{1.4}
\begin{tabular}{@{}cccccc@{}}
$x$ & 0 & 1 & 2 & 3 & 4 \\
\hline
$P(X=x)$ & 0.05 & 0.20 & 0.35 & 0.25 & 0.15 \\
\end{tabular}
\end{center}

\begin{enumerate}[leftmargin=1.5em, itemsep=8pt]
  \item Verify this is a valid probability distribution.\\
        \ansline{All values $\ge 0$ and $0.05+0.20+0.35+0.25+0.15 = 1.00$. Valid.}

  \item Find $P(1 \le X \le 3)$. Interpret your answer in context.\\
        \ansline{$P(1 \le X \le 3) = 0.20+0.35+0.25 = 0.80$.}\\[1pt]\ansline{About 80\% of students read between 1 and 3 books over summer break.}

  \item Find $P(X \ge 3)$. Interpret your answer in context.\\
        \ansline{$P(X \ge 3) = 0.25 + 0.15 = 0.40$.}\\[1pt]\ansline{About 40\% of students read 3 or more books over summer break.}

  \item Describe the shape of this distribution. Where is it centered?\\
        \ansline{Single peak at 2; roughly symmetric (slightly heavier toward 3--4).}\\[1pt]\ansline{Centered near 2 books (mean $= 2.25$).}

  \item Write an interpretation of this distribution in context (shape, center, spread).\\
        \ansline{Summer reading is roughly symmetric with a peak at 2 books, centered}\\[1pt]\ansline{near 2 books, and spread from 0 to 4 books; reading 0 is rare (5\%).}
\end{enumerate}
\end{tcolorbox}

\vspace{0.12in}

\begin{tcolorbox}[colback=greenbg!60, colframe=greenacc, title=\textbf{Tier E --- Extension},
    arc=2mm, boxrule=0.8pt, left=4mm, right=4mm, top=3mm, bottom=3mm]
\small
Two classes were surveyed about the number of after-school activities they participate in.

\textbf{Class A:}
\begin{center}
\renewcommand{\arraystretch}{1.3}
\begin{tabular}{@{}ccccc@{}}
$x$ & 0 & 1 & 2 & 3 \\
\hline
$P(X=x)$ & 0.40 & 0.35 & 0.20 & 0.05 \\
\end{tabular}
\end{center}

\textbf{Class B:}
\begin{center}
\renewcommand{\arraystretch}{1.3}
\begin{tabular}{@{}ccccc@{}}
$x$ & 0 & 1 & 2 & 3 \\
\hline
$P(X=x)$ & 0.10 & 0.25 & 0.40 & 0.25 \\
\end{tabular}
\end{center}

\begin{enumerate}[leftmargin=1.5em, itemsep=8pt]
  \item Verify both distributions are valid.

        \noindent\begin{minipage}[t][0.2in][t]{\linewidth}\ans{A: $0.40+0.35+0.20+0.05 = 1$; B: $0.10+0.25+0.40+0.25 = 1$; all values $\ge 0$.}\end{minipage}

  \item Compare the two distributions: which class tends to participate in \emph{more}
        activities? Support your answer using probabilities.\\
        \ansline{Class B. B puts 0.65 on 2--3 activities; A puts 0.75 on 0--1.}\\[1pt]\ansline{So Class B students tend to do more activities.}

  \item For each class, find $P(X \ge 2)$ and compare.

        \noindent\begin{minipage}[t][0.25in][t]{\linewidth}\ans{A: $0.20+0.05 = 0.25$; B: $0.40+0.25 = 0.65$. Class B is much higher.}\end{minipage}

  \item If a student is chosen at random from Class A, what is the probability they
        participate in \emph{at least} as many activities as the most common value in
        Class B? Show your work.\\
        \ansline{Most common value in Class B is $x = 2$ ($P = 0.40$).}\\[1pt]\ansline{Class A: $P(X \ge 2) = 0.20 + 0.05 = 0.25$.}

  \item Describe the difference in shape between the two distributions.\\
        \ansline{Class A is skewed right (peak at 0, tail toward 3).}\\[1pt]\ansline{Class B is skewed left (peak at 2, tail toward 0).}
\end{enumerate}
\end{tcolorbox}

\end{document}
